\section{Chapter 1: First steps}\label{sec:14-appendix-solutions:01-chapter-1:1}

\textbf{1. β-reduce \((\lambda x.\lambda y.\, y\, x)\, a\, b\) and compare to \(K\).}

Application associates to the left, so this is \(((\lambda x.\lambda y.\, y\, x)\, a)\, b\): \[
(\lambda x.\lambda y.\, y\, x)\, a\, b
\;\longrightarrow_\beta\; (\lambda y.\, y\, a)\, b
\;\longrightarrow_\beta\; b\, a
\] Step 1 substitutes \(a\) for \(x\) in \(\lambda y.\, y\, x\), giving \(\lambda y.\, y\, a\). Step 2 substitutes \(b\) for \(y\) in \(y\, a\), giving \(b\, a\).

The result is \textbf{not} simply one of the two inputs, the way \(K\)'s result always is. \(K\,a\,b \to_\beta a\) --- \(K\) \emph{discards} its second argument and returns the first, verbatim. Here, \((\lambda x.\lambda y.\, y\, x)\, a\, b
\to_\beta b\, a\) --- both arguments survive, but with the \emph{second} one applied to the \emph{first}. This term is sometimes called the ``Thrush'' combinator (\(T\), satisfying \(T\, x\, y = y\, x\)): it flips the order of application rather than discarding anything, which is genuinely different behavior from \(K\), not just a relabeling of it.

\textbf{2. \texttt{Vec.toList : Vec α n → List α}}

\begin{lstlisting}
def Vec.toList : Vec α n → List α
  | Vec.nil => []
  | Vec.cons a rest => a :: Vec.toList rest
\end{lstlisting}

\texttt{Vec.replicate}'s type, \texttt{(n : Nat) → Vec α n}, is genuinely dependent: the \emph{return} type \texttt{Vec α n} mentions \texttt{n}, the value just supplied as the argument. \texttt{Vec.toList}'s type, \texttt{Vec α n → List α}, is not --- its return type \texttt{List α} never mentions \texttt{n} at all, even though its \emph{argument} type happens to be dependent (\texttt{Vec α n}, one specific type per length). Taking a dependently-typed \emph{input} does not automatically make a function dependent; what matters is whether the \emph{output} type varies with the input's \emph{value}. \texttt{Vec.toList} throws the length away on the way out, the same way \texttt{Vec α n}'s own length information disappears once converted to a plain \texttt{List α}.

\textbf{3. A second \texttt{Σ n : Nat, Fin n}, and why \texttt{Σ n : Nat, n > 0} fails}

\begin{lstlisting}
def anotherSigma : Σ n : Nat, Fin n := ⟨5, ⟨0, by decide⟩⟩
\end{lstlisting}

(\texttt{0 : Fin 5} is valid since \(0 < 5\); any pair \texttt{⟨k, ⟨i, h⟩⟩} with \texttt{i < k} works.) \texttt{Σ n : Nat, n > 0} fails to type-check because \texttt{n > 0 : Prop} (\texttt{Sort 0}), while Lean's \texttt{Sigma} is declared as \texttt{structure Sigma \{α : Type u\} (β : α → Type v)} --- the second component's \emph{family} must land in some \texttt{Type v} (\texttt{Sort (v+1)} for \texttt{v ≥ 0}), which \texttt{Prop} (\texttt{Sort 0}) is not. \texttt{Fin n : Type} clears that bar; \texttt{n > 0 : Prop} does not. This is exactly why \texttt{∃} exists as \texttt{Sigma}'s \texttt{Prop}-restricted cousin (Chapter 1 §5) rather than everyone just writing \texttt{Σ} everywhere.

\textbf{4. \texttt{Path.append}'s signature as nested Π-types}

\texttt{Path.append : \{u v w : V\} → Path Q u v → Path Q v w → Path Q u w}, treating the implicit binders as outer Π's, unfolds one binder at a time:

\begin{tabular}{llll}
\toprule
Level & \(A\) & \(B(\cdot)\) & Dependent? \\
\midrule
\bottomrule
1 (binds \(u\)) & \(V\) & \(\prod_{v:V}\prod_{w:V} (\mathrm{Path}_Q\,u\,v \to \mathrm{Path}_Q\,v\,w \to \mathrm{Path}_Q\,u\,w)\) & Yes --- mentions \(u\) \\
2 (binds \(v\)) & \(V\) & \(\prod_{w:V} (\mathrm{Path}_Q\,u\,v \to \mathrm{Path}_Q\,v\,w \to \mathrm{Path}_Q\,u\,w)\) & Yes --- mentions \(v\) \\
3 (binds \(w\)) & \(V\) & \(\mathrm{Path}_Q\,u\,v \to \mathrm{Path}_Q\,v\,w \to \mathrm{Path}_Q\,u\,w\) & Yes --- mentions \(w\) \\
4 (binds \(p\)) & \(\mathrm{Path}_Q\,u\,v\) & \(\mathrm{Path}_Q\,v\,w \to \mathrm{Path}_Q\,u\,w\) & No --- constant in \(p\) \\
5 (binds \(q\)) & \(\mathrm{Path}_Q\,v\,w\) & \(\mathrm{Path}_Q\,u\,w\) & No --- constant in \(q\) \\
\end{tabular}


The first three levels are genuinely dependent Π-types: each \(B\) mentions the vertex just bound, since it appears as an index into \texttt{Path} later in the signature. The last two levels are Π-types that happen to collapse to ordinary arrows: once \(u, v, w\) are fixed, \texttt{Path Q u w}'s \emph{type} no longer depends on the specific \emph{proof term} \texttt{p} or \texttt{q} supplied, only on the vertices already bound --- exactly the ``\(B(x)\) not depending on \(x\)'' collapse case from Chapter 1 §3/§5.
